% CLAUDE REPORT::WR-BRIDGE-0001::2026-11-11T02:50-08:00
\documentclass[11pt]{article}
\usepackage{fontspec}
\usepackage[margin=0.9in]{geometry}
\usepackage{amsmath,amssymb,amsthm,booktabs,xurl,hyperref}
\newtheorem{theorem}{Theorem}
\newtheorem{lemma}[theorem]{Lemma}
\newtheorem{corollary}[theorem]{Corollary}
\theoremstyle{definition}
\newtheorem{definition}[theorem]{Definition}
\newcommand{\SAT}{\mathsf{SAT}}
\newcommand{\UNSAT}{\mathsf{UNSAT}}
\newcommand{\WR}{\overline W_R}
\newcommand{\st}[1]{\textbf{[#1]}}
\title{WR-BRIDGE-0001: The Worst-Case Minimum-Derivation Target\\
and the Exact Location of the Missing Bridge}
\author{Claude research agent (scheduled routine; not one of LSPV-1..5)}
\date{11 November 2026, 02:50 America/Los\_Angeles \quad base commit \texttt{9c91ff01}}
\begin{document}\maketitle

\begin{center}\fbox{\parbox{0.94\linewidth}{
\textbf{Bottom line.} The shared objective
$\WR(N)=\max_{F\in\UNSAT,|\langle F\rangle|\le N}\min_{\pi:F\vdash_R\bot}\mathrm{cost}_R(\pi)$
is, for every fixed Cook--Reckhow system $R$ and $\mathrm{cost}=$ size, exactly the
worst-case proof-size function of $R$. Letting the optimal derivation be chosen for free is
\emph{nondeterminism}; so $\WR$ is a coNP-proof-length measure.
\st{PROVED/EXACT} (Thm.~2): for every $R$ with an already-known superpolynomial lower
bound (Resolution, Res($k$), Cutting Planes, Polynomial Calculus, bounded-depth Frege, \dots)
the proposed bridge $\SAT\in P\Rightarrow\WR(N)\le\mathrm{poly}(N)$ is \emph{equivalent to}
$P\ne NP$; it is not a weaker sufficient step.
\st{PROVED/EXACT} (Thm.~3): for an optimal proof system the bridge is a theorem, but then the
lower bound required is equivalent to $NP\ne coNP$.
\st{FALSIFIED/RETIRED} (Thm.~4): the pointwise simulation ``every poly-time SAT
algorithm's run yields a poly-size Resolution derivation'' (Tseitin/Gaussian elimination).
\st{OPEN}: the bridge for Frege / Extended Frege; $P$ versus $NP$ remains OPEN.}}\end{center}

\section{Mapping the user's idea to exact objects}
The multisign DAG hypergraph of a 3-CNF $F$ is the signed clause hypergraph
(vertices = variables, hyperedges = clauses with literal signs) together with the DAG of
derived clauses. ``Minimal reduced logical state'' for $F\in\UNSAT$ is $\bot$ (or, for a
reduction calculus without $\bot$, a fixed canonical contradiction); ``minimal work'' is the
minimum cost of an $R$-derivation DAG. Granting free knowledge of the best sequence is
legitimate for a \emph{lower} bound on executing it: any execution that \emph{is} an
$R$-derivation costs at least $\min_\pi\mathrm{cost}_R(\pi)$. \st{PROVED/EXACT}, trivially.
The only nontrivial question is which algorithms' executions are (or compile to)
$R$-derivations. This is the transfer problem, and it is where all difficulty lives.
The corpus already used the signed (2,2) hypergraph (LSPV-2, SIL/M22 ladder) and the
Cook--Reckhow boundary (LSPV-5/163); this report reconciles the new directive with both.

\begin{definition}
A proof system $R$ is a poly-time $V_R(F,\pi)$ with $F\in\UNSAT\iff\exists\pi\,V_R(F,\pi)=1$
(Cook--Reckhow \cite{cr}). $R$ is \emph{p-bounded} if $\WR(N)\le\mathrm{poly}(N)$.
$\mathrm{Bridge}(R)$ denotes the implication $\SAT\in P\Rightarrow R$ is p-bounded.
\end{definition}

\section{Theorems}
\begin{theorem}[Cook--Reckhow, \st{KNOWN THEOREM}]
Some proof system is p-bounded iff $NP=coNP$.
\end{theorem}
\noindent Warning recorded for the red team: a literature-search summary obtained this run
asserted that $P=NP$ ``would make every proof system, Extended Frege included, polynomially
bounded.'' \st{FALSIFIED}: $P=NP$ gives \emph{some} p-bounded system (e.g.\ empty proofs
checked by the decider); Resolution would still need $2^{\Omega(n)}$-size refutations of
$\mathrm{PHP}^{n+1}_n$ \cite{haken}, a theorem independent of $P$ vs $NP$.

\begin{theorem}[Bridge collapse for weak systems, \st{PROVED/EXACT}]
If $R$ has an unconditional superpolynomial lower bound $\WR(N)\ne N^{O(1)}$, then
$\mathrm{Bridge}(R)\iff P\ne NP$.
\end{theorem}
\begin{proof}
($\Rightarrow$) Contrapositive of the bridge with the known lower bound gives $\SAT\notin P$.
($\Leftarrow$) If $P\ne NP$ the antecedent $\SAT\in P$ is false, so the implication holds vacuously.
\end{proof}
\noindent Consequence: for Resolution, Res($k$), Cutting Planes \cite{pud}, Polynomial
Calculus, AC$^0$-Frege, proving $\mathrm{Bridge}(R)$ is \emph{exactly as hard as} $P\ne NP$.
The ``weakest sufficient bridge'' is not weaker than the goal for any system whose lower
bound is already in hand. Any proposed proof of such a bridge must therefore be audited as
a would-be proof of $P\ne NP$ (relativization: the statement is a non-relativizing
consequence only through the antecedent; natural-proofs/algebrization audits apply to the
bridge argument itself).

\begin{theorem}[Optimal systems, \st{PROVED/EXACT}]
If $R$ is optimal (every system $Q$ has $\mathrm{size}_R(F)\le\mathrm{poly}(\mathrm{size}_Q(F),|F|)$),
then $\mathrm{Bridge}(R)$ holds, and $\WR$ superpolynomial $\iff NP\ne coNP$.
\end{theorem}
\begin{proof}
$P=NP\Rightarrow NP=coNP\Rightarrow$ some p-bounded $Q$ (Thm.~1) $\Rightarrow R$ p-bounded by
optimality. The equivalence is Thm.~1 plus optimality.
\end{proof}
\noindent Existence of optimal systems is \st{OPEN} (Krajíček--Pudlák \cite{kp}: implied by
$NE=coNE$). So the bridge strategy has a dichotomy: weak $R$ $\Rightarrow$ bridge $=$ whole
problem; optimal $R$ $\Rightarrow$ lower bound $=$ $NP\ne coNP$ (believed harder than $P\ne NP$).
The only genuinely intermediate regime is Frege/Extended Frege, where neither half is known.
There, Pich--Santhanam \cite{ps} prove $\mathrm{Bridge}(EF)$-type statements \emph{under
additional hypotheses} (conditional), and Cook's correspondence \cite{cook75} gives
$\mathrm{Bridge}(EF)$ if $P=NP$ is witnessed by an algorithm whose correctness is provable
in $PV$. \st{CONDITIONAL}. No EF lower bound is known; Li's 2025 conjecture \cite{li}
concerns unprovability of $NP=coNP$ in Cook's theory. Best next target: state precisely
which weakening of the PS hypotheses suffices for $\mathrm{Bridge}(EF)$.

\begin{theorem}[Pointwise simulation fails for Resolution, \st{FALSIFIED/RETIRED} as a transfer]
There are UNSAT 3-CNFs (Tseitin formulas on bounded-degree expanders) refuted by a
deterministic $O(n^3)$ algorithm (Gaussian elimination over $\mathbb F_2$) while every
Resolution refutation has size $2^{\Omega(n)}$ (Urquhart \cite{urq}, Ben-Sasson--Wigderson
\cite{bsw}). Hence ``execution of any poly-time algorithm compiles to a poly-size
Resolution DAG'' is false; the worst-case/pointwise simulation theorem cannot be proved for
$R$ = Resolution, and it cannot be proved for any $R$ with a lower bound on a family in
$P$-decidable UNSAT (PHP for Polynomial Calculus vs.\ Cutting Planes, etc.).
\end{theorem}
\noindent Thus the ``fixed combined inference system'' must at least p-simulate
Resolution + $\mathbb F_2$-linear reasoning + Cutting Planes; but any such combination
with a proved lower bound again falls under Thm.~2. Adaptive strategy switching does not
escape: a switching algorithm is still a single poly-time algorithm, and its run on the
hard family must still be an $R$-derivation.

\section{Computational sanity check (\st{COMPUTATIONAL})}
Script \path{research/agents/claude/wr-worstcase/tree_res_exact.py} computes the exact
minimum tree-like Resolution refutation size (= minimum decision-tree nodes) by exhaustive
memoized DP over partial assignments, and the XOR-row operations of Gaussian elimination.
Raw output: \path{php_results.json}, \path{tseitin_results.json}. Two chunks, each
$<6$\,s, per the 30-second rule. Deterministic, no seeds.
\begin{center}\begin{tabular}{lrrrr}\toprule
instance & vars & clauses & min tree-Res size & Gauss XOR ops\\\midrule
PHP$_{2\to1}$ & 2 & 3 & 5 & --\\
PHP$_{3\to2}$ & 6 & 9 & 21 & --\\
PHP$_{4\to3}$ & 12 & 22 & 85 & --\\
Tseitin $K_4$ & 6 & 16 & 55 & 5\\
Tseitin $K_{3,3}$ & 9 & 24 & 151 & 11\\
Tseitin $Q_3$ & 12 & 32 & 255 & 20\\\bottomrule\end{tabular}\end{center}
PHP sizes fit $(4^{m+1}-1)/3$ for $m\le3$; Tseitin $Q_3$ equals $2^{8}-1$. These finite values
are illustrations of the Thm.~4 gap (exponential-shaped tree-Res vs.\ linear Gauss work),
not asymptotic evidence. Hostile control: Gaussian elimination is a correct poly-time
algorithm that is not a Resolution derivation, which is exactly the counterexample class.

\section*{CORPUS REVIEWED}
\path{AGENTS.md}, \path{CODEX.md}, root \path{main.tex} status lines (P vs NP OPEN at
lines 31, 83, 112), LSPV-5/163 PROOF-SYSTEM-FALLBACK (Cook--Reckhow boundary, unknown-$K$
firewall), LSPV-1/412 RUP-native-resolution bridge, Claude world ledger (row classification),
coordination handoffs of 2026-11-11, and repository search for ``Cook--Reckhow'',
``p-bounded'', ``worst-case''.
\section*{DUPLICATION / CONTRADICTION AUDIT}
Thm.~1 duplicates LSPV-5/163 (consistent). Thms.~2--4 as applied to the new $\WR$ directive
are not stated in the corpus. No contradiction found. The directive's phrase ``weakest
sufficient bridge'' is overstated for weak $R$ (Thm.~2).
\section*{WHAT CHANGED RELATIVE TO ROOT MAIN}
Nothing in root \path{main.tex} changes; this agent does not own a region. LSPV-5 may adopt
Thms.~2--4 into the barrier ledger.
\section*{ADOPTED / REJECTED / RETIRED PRIOR COMPONENTS}
Adopted: Cook--Reckhow boundary. Retired: pointwise algorithm-to-Resolution simulation.
\section*{STALE OR OVERSTATED CLAIMS FOUND}
The external summary claim ``$P=NP$ makes every proof system p-bounded'' (false; see Thm.~1 note).
\section*{Post-quantum ledger}
Collision/lossiness (item 4) and reduction architecture (item 5): ANALOGY/CANDIDATE only;
no exact transfer to $\mathrm{Bridge}(R)$ found.

\begin{thebibliography}{9}
\bibitem{cr} S.~Cook, R.~Reckhow, The relative efficiency of propositional proof systems, JSL 44 (1979).
\bibitem{haken} A.~Haken, The intractability of resolution, TCS 39 (1985).
\bibitem{urq} A.~Urquhart, Hard examples for resolution, JACM 34 (1987).
\bibitem{bsw} E.~Ben-Sasson, A.~Wigderson, Short proofs are narrow, JACM 48 (2001).
\bibitem{pud} P.~Pudlák, Lower bounds for resolution and cutting plane proofs and monotone computations, JSL 62 (1997).
\bibitem{kp} J.~Krajíček, P.~Pudlák, Propositional proof systems, the consistency of first order theories and the complexity of computations, JSL 54 (1989).
\bibitem{cook75} S.~Cook, Feasibly constructive proofs and the propositional calculus, STOC 1975.
\bibitem{ps} J.~Pich, R.~Santhanam, Towards $P\ne NP$ from Extended Frege lower bounds, arXiv:2312.08163.
\bibitem{li} J.~Li, ECCC TR25-187 (2025), \url{https://eccc.weizmann.ac.il/report/2025/187/}.
\end{thebibliography}
\end{document}
% END_OF_VERIFIED_REPORT::CLAUDE::WR-BRIDGE-0001
